% !PRINTED 230
\weberpage{230}
%<<CONT
of order 60 on five letters. This permutation group is obtained by right-multiplying the right cosets
\[
Q, \quad Q\Theta, \quad Q\Theta^2, \quad Q\Theta^3, \quad Q\Theta^4
\]
by all the elements $\sigma$ of the icosahedral group, that is, by forming
\[
Q\sigma, \quad Q\Theta\sigma, \quad Q\Theta^2\sigma, \quad Q\Theta^3\sigma, \quad Q\Theta^4\sigma
\]
and examining the resulting permutation of these right cosets. For the generating substitutions of the icosahedral group $\sigma = \Theta, \psi, \chi$ we thus obtain the permutations
\[
(1, \ \Theta, \ \Theta^2, \ \Theta^3, \ \Theta^4) \qquad \sigma = \Theta
\]
\[
\begin{pmatrix} Q, & Q\Theta, & Q\Theta^2, & Q\Theta^3, & Q\Theta^4 \\ Q, & Q\Theta^4, & Q\Theta^3, & Q\Theta^2, & Q\Theta \end{pmatrix} \qquad \sigma = \psi
\]
\[
\begin{pmatrix} Q, & Q\Theta, & Q\Theta^2, & Q\Theta^3, & Q\Theta^4 \\ Q, & Q\Theta^3, & Q\Theta^4, & Q\Theta, & Q\Theta^2 \end{pmatrix} \qquad \sigma = \chi.
\]

The last is found from the formulae (1) and (2), according to which
\[
\Theta\chi = \varphi\Theta^3, \quad \Theta^2\chi = \chi\varphi^2\Theta^4, \quad \Theta^3\chi = \varphi^2\Theta, \quad \Theta^4\chi = \varphi\chi\Theta^2.
\]

One sees that these permutations all belong to the first kind, and that therefore the permutation group in question can be no other than the alternating group on five letters (Bd. I, \S.\,177). From this it also follows that the icosahedral group is simple, and consequently coincides with the group which we already designated provisionally as the icosahedral group in \S.\,34.

\begin{center}
\S.\,60.\\[2pt]
\gesperrt{The Fundamental Forms of the Icosahedral Group.}
\end{center}

We have found above one of the three fundamental forms of the icosahedral group, $f$:
\weq{1}{f = x_1 x_2 (x_1^{10} + 11 x_1^5 x_2^5 - x_2^{10}).}
We therefore still lack two of these forms, one of the 20th and one of the 30th degree.

The fundamental form of degree 20 is obtained as the \gesperrt{Hessian} covariant of $f$. We set it
\weq{2}{H = \frac{1}{121}\,[f''(x_1,x_1)\, f''(x_2,x_2) - f''(x_1,x_2)^2],}
%CONT>>
